\chapter{Introdução}

A introdução é um capítulo fundamental do trabalho acadêmico, pois apresenta o tema, sua delimitação e a problemática que orienta a pesquisa. Nesse sentido, é necessário contextualizar o assunto, evidenciando sua relevância para a área de estudo e justificando a importância da investigação proposta.

Inicialmente, apresenta-se o tema e o contexto da área de pesquisa, considerando que o leitor pode não estar familiarizado com o assunto. Em seguida, realiza-se a delimitação do tema, especificando o recorte adotado, as tecnologias ou abordagens utilizadas e o público ou área de aplicação.

Posteriormente, define-se a problemática da pesquisa, explicitando o problema que se pretende investigar ou resolver. Essa etapa é essencial para demonstrar a motivação do trabalho e sua contribuição potencial.

O presente trabalho tem como objetivo geral detalhar o modelo de escrita para a produção do Trabalho de Conclusão de Curso. Como objetivos específicos, busca-se realizar a revisão bibliográfica sobre o tema, criar um tutorial sobre a utilização do Overleaf e aplicar os conhecimentos adquiridos na edição de documentos acadêmicos.

O capítulo pode ser escrito em texto corrido como vimos acima, ou pode ser estruturado em seções, como apresentado a seguir:

\section{Contextualização}

A introdução é um capítulo fundamental do trabalho acadêmico, pois apresenta o tema, sua delimitação e a problemática que orienta a pesquisa. Nesse sentido, é necessário contextualizar o assunto, evidenciando sua relevância para a área de estudo e justificando a importância da investigação proposta.

Inicialmente, apresenta-se o tema e o contexto da área de pesquisa, considerando que o leitor pode não estar familiarizado com o assunto. Em seguida, realiza-se a delimitação do tema, especificando o recorte adotado, as tecnologias ou abordagens utilizadas e o público ou área de aplicação.

\section{Justificativa}

A relevância deste trabalho está associada à necessidade de compreender e estruturar adequadamente a produção de trabalhos acadêmicos, especialmente no contexto do Trabalho de Conclusão de Curso. A correta organização e apresentação das informações contribuem para a clareza, qualidade e rigor científico do estudo.

\section{Problema de Pesquisa}

Posteriormente, define-se a problemática da pesquisa, explicitando o problema que se pretende investigar ou resolver. Essa etapa é essencial para demonstrar a motivação do trabalho e sua contribuição potencial.

\section{Objetivos}

O presente trabalho tem como objetivo geral detalhar o modelo de escrita para a produção do Trabalho de Conclusão de Curso.

Como objetivos específicos, busca-se:

\begin{itemize}
    \item Realizar a revisão bibliográfica sobre o tema;
    \item Criar um tutorial sobre a utilização do Overleaf;
    \item Aplicar os conhecimentos adquiridos na edição de documentos acadêmicos.
\end{itemize}

\section{Metodologia}

A metodologia adotada neste trabalho baseia-se em revisão bibliográfica e estudo prático das ferramentas utilizadas para elaboração de documentos acadêmicos. Além disso, são aplicadas técnicas de organização textual e formatação conforme normas vigentes.

\section{Organização do Trabalho}

Este trabalho está estruturado da seguinte forma: o Capítulo 1 apresenta a introdução, contextualizando o tema e os objetivos da pesquisa. O Capítulo 2 aborda o referencial teórico relacionado ao tema. O Capítulo 3 descreve a metodologia utilizada. O Capítulo 4 apresenta os resultados obtidos. Por fim, são apresentadas as considerações finais do trabalho.
